Coastal storm surges are among the more serious natural hazards for coastal communities. High population density, the importance of coastal cities for maritime activities such as ports and infrastructure mean that these areas are particular vulnerable to flooding, material damage and loss of human life \citep{kaniewski_solar_2016}. 
Along with comprehensive economic and health related consequences, storm surges cause major disruptions in transportation and energy-infrastructure which strains emergency- and healthcare systems for coastal communities \citep{lane_health_2013}. 
It is expected that the economic consequences of storm surges globally will rise to around \$65bn in the year 2050, together with the expenses for building new protective dikes is expected to increase with increasing populations and city expansion along the coastlines \citep{hallegatte_future_2013, prahl_damage_2018}.\\
As global sea level rises, an increase in intensity and frequency of storm surges along coastline is likely to be expected \citep{muis_global_2023}. This in return will amplify the necessity of both proactive and adaptive coastal protection measures \citep{marsooli_climate_2019, dedekorkut-howes_when_2020}.\\

It is therefore essential to develop and maintain advanced modelling tools that can predict and assess storm surges' impact on our coastal societies. Early implementations of storm surge models were based on empirical measurements of pressure differences before, under, and after the occurrence of a storm. This led to the development of empirical formulas that could predict water levels using local support data \citep{massey_history_2007}. Later models were developed and tested by reconstructing historical storm surges and in return created extensive databases with training data used in trend analysis to predict changes in storm surge frequency patterns \citep{dang_dataset_2024, tadesse_database_2021}.\\
Towards the end of the 20th century alongside increases in computational power made it possible to integrate more parameters and make the models more dynamic. This development has made it possible to predict the impacts of storm surges on large scales with high speed \citep{tadesse_database_2021, massey_history_2007}.\\
Today, advanced hydrodynamic models such as the Danish Institute of Hydraulics' (DHI) MIKE21/3 \citep{dhi_mike_2024} and University of North Carolina's ADCIRC are able to simulate complex hydrodynamic environmental states in both 2- and 3D by using local and regional climatological data (wind speed, wind direction, coastal morphology, bathymetry, and geology).\\
Another example is a model capable of handling complex situations is the Super Fast Inundation of Coasts (SFINCS) first described by \cite{leijnse_modeling_2021} which is able to compute flooding from both coastal storm surges and heavy rainfall at the same time.\\
Since most of these models are complex and operate on regional or national scales, low-resolution data is often used to enable real-time simulations in response to upcoming storm surge events. \\

It is therefore equally important to develop models that can quickly and easily provide an overview of potential storm surge scenarios and their impact on a more local scale \citep{balstrom_kirby_inundation}. A solution to this, is raster based inundation models that work well in GIS-environments and are suitable for fast visualisation and future analysis of storm surge induced flooding. Methods such as the bathtub-method, using a simple slicing of elevation values in the elevation model to a certain level \citep{poulter_raster_2008} to cost-distance algorithms, that ensures logical movement of water through the terrain while taking natural barriers into account \citep{williams_geographic_2022, li_delineating_2014}, are well-documented.\\ Raster based inundation models could therefore increase the accessibility of semi-advanced tools for initial storm surge impact assessments. As mentioned by \cite{balstrom_kirby_inundation} and \cite{li_delineating_2014} these models would take advantage of better technology and higher resolution data while keeping the models complexity to a minimum. One of those is the Inundation Model and toolset first published by \cite{balstrom_kirby_inundation}.\\

The Danish Ministry of Industry, Business and Financial Affairs \citep{erhvervsministeriet_fremtidige_nodate} notes that there have been registered 27 storm surges in the in Denmark from 1991 through 2017. This accounts to approximately one every year. The ministry expects this to rise to two yearly occurrences in 2100. The last official recognised storm surge that hit Denmark occurred in mid October of 2023. In an effort to strengthen the Danish municipalities emergency response efforts, it is essential to critically evaluate the methods and models used in storm surge analysis. This includes verifying the models capability to correctly replicate past storm surge events, which serves as validation of the models ability to simulate the extent of future storm surge events.\\ 